% Lösungen Einheit 2 von 3 – Produkte aus Polynom und e-Funktion (zusammenbau v0.1)
\einheitenkopf{Einheit 2 von 3 · Produkte aus Polynom und e-Funktion}

\erg{13}{a) für $x \to -\infty$ \quad b) für $x \to +\infty$: $f(x) \to +\infty$; für $x \to -\infty$: $f(x) \to 0$ \quad c) für $x \to +\infty$: $f(x) \to -\infty$; für $x \to -\infty$: $f(x) \to 0$, Annäherung von oben \quad d) $f(x) \to 0$: der e-Faktor fällt schneller, als $x + 4$ wächst; der Graph nähert sich von oben der $x$-Achse. \quad e) für $x \to +\infty$: $f(x) \to +\infty$; für $x \to -\infty$: $f(x) \to 0$ \quad f) $w(t) \to 40$: das Produkt geht gegen null, der Summand bleibt; der Pegel sinkt auf Dauer wieder auf 40 cm. \quad g) Nullstelle $x = 5$ (der e-Faktor ist nie null); für $x \to +\infty$: $f(x) \to +\infty$; für $x \to -\infty$: $f(x) \to 0$, von unten \quad h) für $x \to -\infty$: $f_a(x) \to 0$; für $x \to +\infty$: $f_a(x) \to -\infty$ – beides für jedes $a$, weil für große $|x|$ der Summand $-x$ das Vorzeichen des Faktors bestimmt}
\erg{14}{a) Fehler: der e-Faktor geht schneller gegen null, als das Polynom wächst. Richtig: $f(x) \to 0$ für $x \to +\infty$. \quad b) $e^{x}$ ist für jedes $x$ positiv, nie null; ein Produkt ist nur null, wenn ein Faktor null ist – also nur für $x - 1 = 0$. \quad c) $c(t) \to 0$: das Medikament wird auf Dauer vollständig abgebaut. \quad d) $(x - 1) \cdot e^{x}$}
